
\documentclass[10pt]{article}

%% The amssymb package provides various useful mathematical symbols
\usepackage{amssymb}
\usepackage{cancel}
\usepackage{ulem} 
\usepackage{authblk}

\begin{document}


 

\author{Alejandro Rivero \thanks{\texttt{arivero@unizar.es}}}
\affil{BIFI - Universidad de Zaragoza, \\ Mariano Esquillor, Edificio I + D \\ 50018 Zaragoza (Spain)}



 

\title{Bootstrapping generations}
 
 \maketitle 
\begin{abstract}
A supersymmetric version of Chew's old democratic bootstrap argument
predicts the existence of three generations of quarks and leptons with one
quark, of type "up", more massive that the other five. 
\end{abstract}


%\begin{keyword}
%% keywords here, in the form: keyword \sep keyword
%example \sep \LaTeX \sep template
%% MSC codes here, in the form: \MSC code \sep code
%% or \MSC[2008] code \sep code (2000 is the default)
%\end{keyword}
 

The origin of the now popular "dual models" can be traced to advocacy,
by Geoffrey Chew in the early sixties, of the study of the S-matrix, 
eventually looking for an unique consistent solution that "bootstraps" 
itself. A guiding principle was "nuclear democracy"; no particle
should be claimed to be elementary. Each particle was 
providing the force needed to bind the others \cite{Chew:1963zza,PhysRevLett.7.394}
and from the view of S-Matrix polologists, each particle
was a composite of others. 

At the start of the seventies, it was noticed that dual models 
with fermions \cite{Ramond:1971gb, Neveu:1971iv} seem to require
a special symmetry, supersymmetry, relating bosonic and
fermionic degrees of freedom. Schwarz \cite{Schwarz:1972hg} has proposed to interpret this need as a "quark-gluon dual model". 

Given the stringy character
of these models, it seems reasonable to consider that the
gluon part, bosonic, is constituted by a string terminated in
a pair of quarks. With this imaginery, we can try a 
version of the democratic bootstrap. We require that scalar
leptons and scalar quarks are
composed of pairs of quarks, bound by some gluonic string.
Still, we keep agnostic about if quarks and leptons
are composed of themselves or other particles, or if
they are truly elementary.

Each possible pair of "string terminations"
will produce one kind of scalar, as a triplet of
colour in the case of pairs quark-quark, and as
a colour singlet in the case of quark-antiquark; i.e.,
we are taking the triplet in $3 \times 3 = 3 + 6$ and the 
singlet in $3 \times \bar 3 = 1 + 8$ when combining 
the colour charges of the constituents. As for the
electric charge, it is simply the sum of the charges
of the two elements of the pair.

Following an emergent tradition, we will call squarks
and sleptons to the scalar particles, and we will refer
to the requirement of matching between the number of
composite scalar particles and the number of degrees of
freedom of the standard model fermions as the "sBootstrap".
This matching, of course, must be fulfilled in
each charge sector, for each value of electric 
and colour charge, but we do not need to consider
the later in the calculations, as simply we have
triplets coinciding with triplets and singlets
with singlets.


So, lets first taste the sBootstrap condition at cask strength:
$N$ generations of
quarks of types $u, d$ should produce $N^2$ scalars of charge
$+1/3$, but we have $2 N$ fermionic degrees of freedom
with such charge, so $N=2$ for the matching in the
"down antiquark" sector. And they also produce, from pairs
of two down-type quarks, a total of $N (N+1)/2$ scalars
of charge $-2/3$, again to be matched to $2 N$ degrees of
freedom in the "up antiquark" sector. So $N=3$ for the
matching in this sector. The mechanism fails, which
is no surprise.

We need thus to dilute our condition for it to be palatable.
We do it in the following simple way: requesting that not every
quark can terminate the string. An argument for this is
to imagine that some of the families have very high
masses, perhaps of the order of the electroweak scale,
so that on one hand they can not reach the relativistic
speed at the extremes of the string, and on the other hand
they would even disintegrate and decay before 
hadronization, not being able to form a durable bound state.

We postulate thus that only some subset of "light quarks"
are in the terminations of the string and then able to
form composites, and we will calculate the size
of this subset. Given that we could even have
"mixed" generations, with one light and a heavy quark,
we have now at our disposal three parameters: the
number of generations $N$, the number of light quarks of
the down kind, say $r$, and the number of light quarks of
the up kind, say $s$. The matching of scalar and
fermionic degrees of freedom is now

\begin{eqnarray}
r s = 2 N \\
{r (r+1) \over 2} = 2N  
\end{eqnarray}

respectively for $+1/3$ and $-2/3$ particles. The integer solution of the system is that $N$ must be half
of an hexagonal number % n (2n-1)

\begin{equation}
2 N = \cancel { \, 1}, 6, \cancel {15}, 28, \cancel {45}, 66, \cancel {91}, 120, ...
\end{equation}

So the smallest admissible solution is $N=3$. We could
consider the requirement of
asymptotic freedom  in the beta function 
of QCD \cite{Gross:1973id,Gross:1973ju} to put an upper  limit to the number of 
solutions, asking the number of flavours $n_f$ to be such that $2 n_f < 33$.
If this limitation applies to all the flavours,
light or heavy, then the solution is unique.

\emph{We obtain N=3 generations of quarks with r=3 light
down quarks but only s=2 light up quarks.}

Perhaps it could be argued that the beta function 
should be built only with the light quarks. In this case
we have still a very short list of solutions

% r2 + r - 4N =0 %
% r =  (+ sqrt( 1 + 16N) -1 )/2
 
\begin{center}
\begin{tabular}{c|cc|c}
N  & s & r & Total light flavours \\
\hline
3      &     2        &      3      &       5 \\
14      &    4       &       7       &     11 \\
33      &    6      &        11       &    17 
%60          8              15
\end{tabular}
\end{center}

Where we have included the third solution only because 17
flavours are very near of the theoretical bound of 16.5. 
Of course all the solutions beyond $N=3$ already contain
a lot of heavy generations. Note that we always have $r = 2 s -1$.

To search for another source of uniqueness, lets look
the slepton sector. 

For the charged leptons, we have nothing new. Every
charged scalar lepton is a composite of a quark and
antiquark of different type, so $ r s = 2 N$. 

For neutral leptons, we have some options.
On the Standard Model side, we could consider the
known degrees of freedom of neutrinos, only as left-
handed particles, or add a right-handed neutrino in
each generation.
On the composition side, the question is what group
must be use to classify the colour neutral composites. 
Given that
they are similar to mesons (they could even be interpreted
as the mesons themselves) we will classify
them with $SU(r+s)$. We could try $U(r+s)$, having
an extra neutral state, but then we had a odd number
of neutral states and no solution.

So we have finally two possible equations.

1) With both left and right handed neutrinos:
\begin{equation}
r^2 + s^2 -1 = 4N
\end{equation}


2) With only left handed neutrinos
\begin{equation}
r^2 + s^2 -1 = 2N
\end{equation}

And we can discard the second option: when combined with
the previous equations, it does not produce an
integer solution. Interestingly, the sBootstrap can be extended to
the lepton sector only if we use the extra degrees
of freedom that come with the addition of right-handed
neutrinos.

Furthermore, the extension produces a system of
equations with an unique integer solution,
the one with $N=3$.

So, even if we do not want to use QCD beta-function
as an argument, we also have an unique solution if
we ask the the sBootstrap to produce all the scalars
in the SSM. 

To conclude: we
have found that the sBootstrap requisites predict
that the number of fermion families must be \emph{three}.

Looking the solution group-theoretically, we can say that
the scalars are produced from 3 "down quarks" and 2 "up quarks"
according the decomposition of a flavour group 
$SU(5)$ to $SU(3) \times SU(2)$

We get the sleptons extracted from a {\bf 24} of this
flavour group

\begin{equation}
 24=(1,1)+(3,1)+(2,3)+(2, \bar 3)+(1,8)
\end{equation}

And the squarks from the two
sextets that appear in the
decomposition of {\bf 15}.

\begin{equation}
15=(3,1)+(2,3)+(1,6)
\end{equation}

And similarly the anti-squarks.

In the solution, the number of "light" down-type quarks is $r=3$, equal to the number $N=3$ of generations, but the 
number of "light" up-type quarks is $s=2$.

Having a third generation discovered with one massive "up type" quark, this
should be considered a striking \sout{prediction} postdiction from 
the theory of dual models,
or string theory, how they call it nowadays.


%% The Appendices part is started with the command \appendix;
%% appendix sections are then done as normal sections
%\appendix

%\section{Section in Appendix}
%\label{appendix-sec1}

%% References
%%
%% Following citation commands can be used in the body text:
%% Usage of \cite is as follows:
%%   \cite{key}         ==>>  [#]
%%   \cite[chap. 2]{key} ==>> [#, chap. 2]
%%

%% References with bibTeX database:

\bibliographystyle{plain}

\bibliography{Bootstrap}

\end{document}

%%
%% End of file `elsarticle-template-num.tex'.
